\begingroup
\renewcommand{\arraystretch}{1.45}
\fancyhead[L]{}
\fancyhead[C]{\small Das wunderbare Pentagramm}
\providecommand{\EqOverGt}{\mathrel{\substack{=\\>}}}
\par\medskip\Needspace{5\baselineskip}\pdfbookmark[1]{Art. 5}{pentagramma-art-5}\centerline{[5.]}\nobreak\smallskip
Gleichung der Punkte auf der Kegelfläche, auf der die Punkte $(1)\allowbreak(2)\allowbreak(3)\allowbreak(4)\allowbreak(5)$ liegen; die Spitze des Kegels liegt im Mittelpunkt der Kugel, der zugleich der Koordinatenursprung ist. Die $x$-Achse verläuft durch den Punkt $(3)$, also verläuft die $yz$-Ebene durch $(1)$ und $(5)$, und die $y$-Achse verläuft durch $(1)$
\[
\begin{array}{c@{\qquad}c@{\qquad}c@{\qquad}c} &x&y&z\\[9pt] (3)&1&0&0\\[9pt] (4)&\cos1&0&\sin1\\[9pt] (5)&0&\cos3&\sin3\\[9pt] (1)&0&1&0\\[9pt] (2)&\cos5&\cos4&-\cos3\mathord{.}\sin5 \end{array}
\]
Gleichung
\[
(z\mathord{.}\cos1-x\mathord{.}\sin1)(z\mathord{.}\cos3-y\mathord{.}\sin3)\cos2=xy
\]



\Needspace{24\baselineskip}
\SourceLeaf{496}

oder
\[
zz=xz\surd{\alpha}+yz\surd{\gamma}+\frac{1+\alpha+\gamma}{\surd{\alpha\gamma}}xy
\]
Durch Veränderung der Koordinatenebenen lässt sich dieselbe in die Form
bringen
\[
z'z'=Lx'x'+My'y'
\]
Löst man die Gleichung auf
\[
t(2t-1)^2=\alpha\beta\gamma\delta\varepsilon(t-1)
\]
welche eine negative $(G)$ und zwei positive Wurzeln $(G',G'')$ hat, so wird
\[
\begin{gathered} Gz'z'+G'x'x'+G''y'y'=0\\[9pt] GG'G''=-\tfrac14\alpha\beta\gamma\delta\varepsilon\\[9pt] \qquad(G-1)(G'-1)(G''-1)=-\tfrac14\\[9pt] (2G-1)(2G'-1)(2G''-1)=-\alpha\beta\gamma\delta\varepsilon \end{gathered}
\]

\Needspace{18\baselineskip}
Für obiges Beispiel
\[
t(2t-1)^2=20(t-1)
\]
Wurzeln\quad
$-2{,}1973145,\quad+1{,}06931815,\quad+2{,}1279965$

Setzt man
\[
\alpha\beta\gamma\delta\varepsilon=\omega\quad\text{und}\quad \surd{\frac{(18\omega+1)^2}{(3\omega+1)^{3}}}=\cos3\psi
\]
so wird
\[
t=\tfrac13-\tfrac13\cos\psi\mathord{.}\surd{(3\omega+1)}
\]


\Needspace{44\baselineskip}
Das Verhalten der kubischen Gleichung
\[
\frac{t(2t-1)^2}{t-1}=\alpha\beta\gamma\delta\varepsilon
\]
(welche man am bequemsten mit \textsc{Weidenbach}s Tafel auflöst, wo für
$\frac{1-x}{1+x}=y$
gesucht werden muss $\frac{1}{yxx}=\alpha\beta\gamma\delta\varepsilon$, wonach dann
$2t-1=\frac1x$ wird) übersieht man durch folgende Tafel


\SourceLeaf{497}

\par\nopagebreak% portrait-table start
\begingroup\small\setlength{\arraycolsep}{2pt}\[
\resizebox{.84\linewidth}{!}{$\displaystyle \begin{array}{r@{\quad}r@{\quad}r@{\quad}r@{\quad}r@{\quad}r} t&\alpha\beta\gamma\delta\varepsilon&t&\alpha\beta\gamma\delta\varepsilon&t&\alpha\beta\gamma\delta\varepsilon\\[9pt] \infty&\infty&+1.9&+16.6&+1.0&\infty\\[9pt] +10&+400.9&+1.8&+15.2&\mathord{.}\ \mathord{.}\ \mathord{.}\ \mathord{.}\ \mathord{.}&\text{negativ}\\[9pt] +9&+325.1&+1.7&+14.0&-1.0&\infty\\[9pt] +8&+257.1&+1.6&+12.9&-1.618034&+11.0901699\\[9pt] +7&+197.2&+1.5&+12.0&-2&+16.7\\[9pt] +6&+145.2&+1.4&+11.34&-3&+36.7\\[9pt] +5&+101.2&+1.309017&+11.0901699&-4&+64.8\\[9pt] +4&+65.3&+1.3&+11.09&-5&+100.8\\[9pt] +3&+37.5&+1.2&+11.76&-\infty&\infty\\[9pt] +2&+18&+1.1&+15.84&& \end{array}$}
\]\endgroup\par% portrait-table end



\Needspace{9\baselineskip}
Damit also drei reelle Wurzeln existieren, muss
$\alpha\beta\gamma\delta\varepsilon\EqOverGt11.0901699$
oder $\frac{11}{2}+\frac{\surd{125}}{2}$ gelten; die Grenzwerte für $t$ sind also:
$+1.309017=\frac{3+\surd5}{4}$ und $-1.618034=-\frac{1+\surd5}{2}$.

\clearpage
\endgroup
